\section{Aligned profiles and entire-order addition}
\label{clc:sec:analysis}

For $y>0$ define
\begin{equation}\label{clc:eq:F}
 F_{a,s}(y)=y\Phi(-y,s,a),\qquad
 \ell(x)=\frac{e^x}{1+e^x}.
\end{equation}
Thus $F_{1,s}(y)=-\Li_s(-y)$ and $F_{a,0}(e^x)=\ell(x)$.
For a positive real scale $q$ and $\Rea A>-q$, set
\begin{equation}\label{clc:eq:h}
 h_{q,A,s}(x)=q^{-s}F_{1+A/q,s}(e^{qx}),\qquad
 g_q(p)=\frac{\pi}{q}\csc\frac{\pi p}{q}.
\end{equation}
We call $A$ the affine center. It is $q(a-1)$ in the original Lerch
shift, not $a$ itself.

\begin{lemma}[Weighted entire dependence]\label{clc:lem:weighted}
Let $q>0$, $\Rea A>-q$, and
\begin{equation}\label{clc:eq:strip}
 \max(0,-\Rea A)<c<q.
\end{equation}
Then $s\mapsto e^{-cx}h_{q,A,s}(x)$ is entire with values in
$L^1(\R)\cap L^\infty(\R)$. The assertion holds after any fixed
multiplication by a power of $x$, any fixed order or center derivative,
and locally uniformly in $A$. The profiles are smooth in $x$ and obey
\begin{equation}\label{clc:eq:profile-ladder}
 (\partial_x+A)h_{q,A,s}=h_{q,A,s-1}.
\end{equation}
\end{lemma}
\begin{proof}
For $\Rea s>0$, the Lerch integral and the change of variable from its
Laplace parameter to $qu$ give
\begin{equation}\label{clc:eq:fractional-profile}
 h_{q,A,s}(x)=\frac1{\Gamma(s)}
 \int_0^\infty u^{s-1}e^{-Au}\ell(q(x-u))\dd u.
\end{equation}
Every derivative of $\ell$ is $\ell$ times a polynomial in $\ell$.
Consequently $|\ell^{(j)}(v)|\le C_j\ell(v)$ on the real line.
The function $e^{-cx}\ell(qx)$ is bounded and integrable precisely for
$0<c<q$. Multiplying \eqref{clc:eq:fractional-profile} by $e^{-cx}$ and
writing $v=x-u$ separates the norm estimate into this logistic norm and
\[
 \int_0^\infty u^{\Rea s-1}e^{-(c+\Rea A)u}\dd u.
\]
Both are finite. On compact subsets of the initial parameter domain,
derivatives insert only powers of $u$, $\log u$, and smooth Gamma factors;
these are controlled by slightly larger integrable majorants.

Integration by parts in $u$, initially for $\Rea s>1$, proves
\eqref{clc:eq:profile-ladder}. For an arbitrary compact set of orders,
choose an integer $N$ with $\Rea(s+N)>0$ and define
$h_{q,A,s}=(\partial_x+A)^Nh_{q,A,s+N}$. The integration-by-parts identity
makes these definitions consistent on overlaps. They agree with the
convergent Lerch series for $x<0$ and hence with the specified continuation.
The derivative estimates for $\ell$ prove the norm-holomorphic extension.
Finally, choose neighboring admissible weights $c\pm\varepsilon$.
The bound $|x|^k\le C_{k,\varepsilon}e^{\varepsilon|x|}$ proves every
polynomially weighted assertion. Cauchy's formula in the parameters
supplies all fixed parameter derivatives.
\end{proof}

\begin{theorem}[Scaled bilateral transform]\label{clc:thm:transform}
Under \eqref{clc:eq:strip}, for every $s\in\C$,
\begin{equation}\label{clc:eq:transform}
 \int_\R e^{-px}h_{q,A,s}(x)\dd x
 =g_q(p)(p+A)^{-s},\qquad \Rea p=c.
\end{equation}
The power uses the right-half-plane logarithm of $p+A$.
\end{theorem}
\begin{proof}
For $\Rea s>0$, Fubini is justified by the preceding proof. The transform
of \eqref{clc:eq:fractional-profile} factors into $(p+A)^{-s}$ and
\[
 \int_\R e^{-px}\ell(qx)\dd x
 =\frac1q\int_0^\infty\frac{y^{-p/q}}{1+y}\dd y
 =\frac1q\Gamma(p/q)\Gamma(1-p/q)=g_q(p).
\]
The last step is Gamma reflection \cite{DLMF5}. Both sides are entire in
$s$, so the identity extends to all orders. In particular there is no
remaining $\Gamma(s)$ factor and no remaining $q^s$ factor.
\end{proof}

Let $\bq=(q_1,\ldots,q_r)$, $q_j>0$, and $q_*=\min_jq_j$.
For $r\ge2$ put $x_r=\mu-\sum_{j<r}x_j$ and define
\begin{equation}\label{clc:eq:C-def}
 \CC_{\bq,A}(W;\mu)=
 \int_{\R^{r-1}}\prod_{j=1}^r h_{q_j,A,s_j}(x_j)
           \dd x_1\cdots\dd x_{r-1},\qquad W=\sum_js_j.
\end{equation}
For $r=1$ the notation means evaluation of the single profile at $\mu$.
That the integral depends only on $W$ is part of the following theorem.

\begin{theorem}[Entire-order addition at arbitrary positive scales]
\label{clc:thm:addition}
If $\Rea A>-q_*$, \eqref{clc:eq:C-def} is absolutely convergent for all
independent complex orders, depends only on their sum, and is entire in
that sum. With
$G_{\bq}(p)=\prod_jg_{q_j}(p)$, it equals
\begin{equation}\label{clc:eq:C-contour}
 \CC_{\bq,A}(W;\mu)=\frac1{2\pi i}
 \int_{c-i\infty}^{c+i\infty}
 e^{p\mu}G_{\bq}(p)(p+A)^{-W}\dd p,
 \qquad \max(0,-\Rea A)<c<q_*.
\end{equation}
Every fixed mixed order derivative can be moved inside the ordinary
convolution integral. The contour expression is holomorphic for
$|\Ima\mu|<\pi\sum_jq_j^{-1}$.
\end{theorem}
\begin{proof}
Set $u_j(x)=e^{-cx}h_{q_j,A,s_j}(x)$. On the constraint hyperplane the
weight removed from the product is the constant $e^{c\mu}$. Thus the
absolute integral is at most
$e^{c\mu}\|u_r\|_\infty\prod_{j<r}\|u_j\|_1$.
Lemma~\ref{clc:lem:weighted} gives convergence and locally uniform bounds
for all derivatives. The product of Fourier transforms is
$G_{\bq}(p)(p+A)^{-W}$, by Theorem~\ref{clc:thm:transform}.
Along the vertical line this decays like
$\exp(-\pi|\Ima p|\sum_jq_j^{-1})$ times a fixed power of $1+|p|$
on compact parameter sets. Fourier inversion therefore applies. The
convolution is continuous, so equality holds pointwise. The same decay
proves the claimed complex $\mu$ strip and all holomorphic dependence.
\end{proof}

Write $K_{\bq}(\mu)=\CC_{\bq,A}(0;\mu)$; it is independent of $A$ and
is the convolution of the scaled logistic functions. An equivalent
representation, initially for $\Rea W>0$, is
\begin{equation}\label{clc:eq:base-fractional}
 \CC_{\bq,A}(W;\mu)=\frac1{\Gamma(W)}
 \int_0^\infty u^{W-1}e^{-Au}K_{\bq}(\mu-u)\dd u.
\end{equation}
It follows either by the Dirichlet simplex integral in
\eqref{clc:eq:fractional-profile} or by its bilateral transform. The
entire continuation in $W$, not the improper integral on the right,
is used when $\Rea W\le0$.

Scaling is also exact. For $\lambda>0$,
\begin{equation}\label{clc:eq:scaling}
 \CC_{\lambda\bq,\lambda A}(W;\mu)
 =\lambda^{1-r-W}\CC_{\bq,A}(W;\lambda\mu).
\end{equation}
Indeed $h_{\lambda q,\lambda A,s}(x)=\lambda^{-s}h_{q,A,s}(\lambda x)$,
and the $r-1$ integration variables supply the remaining power.
